\begin{frame}[label=title]
  \centering
  \begin{figure}[H]
    \includegraphics[keepaspectratio, width=0.7\paperwidth, height=.6\paperheight]{Logos/FreeFlower.pdf}
  \end{figure}
  \textbf{Informatik: Datenbanken}\\
  Aggregationsfunktionen, Gruppierung, Verzweigungen\\
\end{frame}

\lstset{style=sql, basicstyle=\ttfamily}

% Zeigt in myexercise/myanswer die Nummer der Aufgabe im Skript statt der Foliennummer.
% \renewcommand wirkt nur innerhalb des aktuellen Frames (Gruppe).
\newcommand{\skriptaufgabe}[1]{\renewcommand{\theSharedCounter}{#1}}

\begin{frame}{Inhalte}
  \begin{itemize}[<+->]
    \item Wiederholung: \gls{SQL}
    \item Aggregatsfunktionen: \lstinline|COUNT|, \lstinline|MAX|, \lstinline|MIN|, \lstinline|SUM|, \lstinline|AVG|
    \item Gruppierung: \lstinline|GROUP BY|, \lstinline|HAVING|
    \item Erste \textit{n} Zeilen: \lstinline|LIMIT|
    \item Verzweigungen: \lstinline|CASE WHEN|
  \end{itemize}
\end{frame}

% Wiederholungsfolien: Lösungen erscheinen erst auf der zweiten Folie (\pause)
\begin{frame}[fragile]{Einfache Selektion: \lstinline|SELECT|}{Wiederholung}
  \skriptaufgabe{2.3}
  \begin{myexercise}
    Geben Sie die Namen und Wohnorte aller Mitglieder aus.
    \pause
    \begin{myanswer}
      \begin{lstlisting}
SELECT name, city
FROM users
\end{lstlisting}
    \end{myanswer}
  \end{myexercise}
\end{frame}

\begin{frame}[fragile]{Selektion mit Filter: \lstinline|WHERE|}{Wiederholung}
  \skriptaufgabe{2.4}
  \begin{myexercise}
    Geben Sie die Namen aller weiblichen Mitglieder aus, die zwischen 150 cm und 155 cm gross sind.
    \pause
    \begin{myanswer}
      \begin{lstlisting}
SELECT name, gender, centimeters
FROM users
WHERE gender='female'
AND centimeters BETWEEN 150 AND 155
\end{lstlisting}
    \end{myanswer}
  \end{myexercise}
\end{frame}

\begin{frame}[fragile]{Eindeutige Selektion: \lstinline|SELECT DISTINCT|}{Wiederholung}
  \small
  \skriptaufgabe{2.8}
  \begin{myexercise}
    Geben Sie jeden Wohnort nur einmal aus.
    \pause
    \begin{myanswer}
      \begin{lstlisting}
SELECT DISTINCT city
FROM users
\end{lstlisting}
    \end{myanswer}
  \end{myexercise}
\end{frame}

\begin{frame}[fragile]{Ungefähre Treffer: \lstinline|LIKE|}{Wiederholung}
  \skriptaufgabe{2.14}
  \begin{myexercise}
    Wählen Sie alle Personen mit dem Namen Naomi aus, die nicht aus Berlin kommen.
    \pause
    \begin{myanswer}
      \begin{lstlisting}
SELECT name, city
FROM users
WHERE name LIKE 'Naomi%'
AND city <> 'Berlin'
\end{lstlisting}
    \end{myanswer}
  \end{myexercise}
\end{frame}

\begin{frame}[fragile]{Sortieren: \lstinline|ORDER BY|}{Wiederholung}
  \skriptaufgabe{2.16}
  \begin{myexercise}
    Geben Sie alle Benutzernamen in sortierter Reihenfolge aus (a $\rightarrow$ z).
    \pause
    \begin{myanswer}
      \begin{lstlisting}
SELECT username
FROM users
ORDER BY username ASC
\end{lstlisting}
    \end{myanswer}
  \end{myexercise}
\end{frame}

\begin{frame}[fragile]{Aggregatsfunktionen: \lstinline|COUNT|, \lstinline|MAX|, \lstinline|MIN|, \lstinline|SUM|, \lstinline|AVG|}
  \skriptaufgabe{2.18}
  \begin{myexercise}
    Geben Sie die Anzahl registrierter Mitglieder in Berlin aus. Die resultierende Spalte soll \enquote{Registrierte Mitglieder in Berlin} heissen.
    \pause
    \begin{myanswer}
      \begin{lstlisting}
SELECT COUNT(*) AS `Registrierte Mitglieder in Berlin`
FROM users
WHERE city='Berlin'
\end{lstlisting}
    \end{myanswer}
  \end{myexercise}
\end{frame}

\begin{frame}{Auftrag}{Aggregatsfunktionen}
  \small
  Moodle
  \begin{itemize}
    \item \aufgabebox{Moodle 2.1--2.3}
    \item \challengebox{Moodle 2.4}
  \end{itemize}
  \textbf{Skript}
  \begin{itemize}
    \item \aufgabebox{Aufgaben 2.18--2.21}
  \end{itemize}
\end{frame}

\begin{frame}[fragile]{Filtern nach Gruppieren: \lstinline|HAVING|}
  \small
  \skriptaufgabe{2.22}
  \begin{mychallenge}
    Geben Sie die durchschnittliche Körpergrösse aller Mitglieder in jeder Stadt aus. Zeigen Sie nur die Städte, in denen die Menschen im Durchschnitt zwischen 150 cm und 155 cm gross sind.
    \pause
    \begin{myanswer}
      \begin{lstlisting}
SELECT city, AVG(centimeters) AS `Durchschnittliche Körpergrösse`
FROM users
GROUP BY city
HAVING `Durchschnittliche Körpergrösse` BETWEEN 150 AND 155
\end{lstlisting}
    \end{myanswer}
  \end{mychallenge}
\end{frame}

\begin{frame}[fragile]{Erste \textit{n} Zeilen: \lstinline|LIMIT|}
  \skriptaufgabe{2.26}
  \begin{myexercise}
    Zeigen Sie die Namen und Körpergrösse der 5 grössten Mitglieder an.
    \pause
    \begin{myanswer}
      \begin{lstlisting}
SELECT name, centimeters
FROM users
ORDER BY centimeters DESC
LIMIT 5
\end{lstlisting}
    \end{myanswer}
  \end{myexercise}
\end{frame}

\begin{frame}{Auftrag}{Gruppieren und Begrenzen}
  \textbf{Skript}
  \begin{itemize}
    \item \challengebox{Aufgaben 2.22, 2.23 (\texttt{HAVING})}
    \item \aufgabebox{Aufgaben 2.24--2.26 (\texttt{LIMIT})}
    \item \challengebox{Aufgabe 2.27}
    \item \aufgabebox{Aufgabe 2.28 (\texttt{GROUP BY}, Abgabe auf Moodle)}
  \end{itemize}
\end{frame}

\begin{frame}[fragile]{Verzweigungen: \lstinline|CASE WHEN|}
  \begin{lstlisting}
SELECT name, centimeters,
CASE
    WHEN centimeters > 179 THEN 'Gross'
    WHEN centimeters < 179 THEN 'Klein'
    WHEN centimeters = 179 THEN 'Genau im Durchschnitt'
    ELSE 'unbekannt'
END AS Durchschnittlich
FROM users
WHERE gender='male'
\end{lstlisting}
  Mitglieder ohne Grössenangabe (\lstinline|NULL|) landen im \lstinline|ELSE|-Zweig.
\end{frame}

\begin{frame}{Auftrag}{Verzweigungen}
  \textbf{Skript}
  \begin{itemize}
    \item \readbox{Abschnitte 2.11 (\texttt{LENGTH}) und 2.12 (\texttt{CASE WHEN})}
    \item \aufgabebox{Aufgabe 2.29}
  \end{itemize}
\end{frame}

\begin{frame}{Werbungs-Teil von InstaHub}{Tabellen \texttt{ads} und \texttt{analytics}}
  \begin{figure}[H]
    \centering
    \adjustbox{max width=\textwidth, max totalheight=.75\textheight}{% Tabellenkarten für den Werbe-Teil von InstaHub (Spalten gemäss den InstaHub-Migrationen,
% https://github.com/wi-wissen/InstaHub, database/migrations/create/ads und .../analytics).
% Jede Tabelle ist eine einspaltige Matrix: Zeile 1 = Tabellenname, danach die Kolonnen.
% Die Pfeile zeigen, auf welche id die Kolonnen user_id und photo_id verweisen.
\begin{tikzpicture}[
    karte/.style={
        matrix of nodes, draw, thick, inner sep=0pt, row sep=0pt,
        nodes={anchor=west, font=\scriptsize\ttfamily, inner sep=1.5pt, text height=1.6ex, text depth=.4ex, minimum width=26mm},
        row 1/.style={nodes={font=\scriptsize\bfseries\ttfamily, fill=RoyalBlue!20}},
      },
    verweis/.style={->, >=stealth, thick, RoyalBlue},
  ]
  \matrix (ads) [karte] {
    ads \\ id \\ priority \\ name \\ type \\ url \\ img \\ query \\ created\_at \\ updated\_at \\
  };
  \matrix (analytics) [karte, right=8mm of ads.north east, anchor=north west] {
    analytics \\ id \\ ip \\ device \\ brand\_family \\ brand\_model \\ browser\_family \\ browser\_version \\
    platform\_family \\ platform\_version \\ user\_id \\ photo\_id \\ created\_at \\ updated\_at \\
  };
  \matrix (users) [karte, right=14mm of analytics.north east, anchor=north west] {
    users \\ id \\ username \\ \ldots \\
  };
  \matrix (photos) [karte, right=14mm of analytics.south east, anchor=south west] {
    photos \\ id \\ user\_id \\ description \\ \ldots \\
  };
  % Verweise analytics.user_id -> users.id und analytics.photo_id -> photos.id
  \draw[verweis] (analytics-11-1.east) -- ++(6mm,0) |- (users-2-1.west);
  \draw[verweis] (analytics-12-1.east) -- ++(9mm,0) |- (photos-2-1.west);
\end{tikzpicture}
}
  \end{figure}
\end{frame}

\begin{frame}[fragile]{Attribute von der Tabelle \texttt{analytics}}{Was wird alles erfasst beim Browsen?}
  \scriptsize
  \begin{table}[H]
    \centering
    \begin{tabular}{rl}
      \textbf{Feld}                 & \textbf{Beschreibung}                   \\\midrule
      \commandbox|id|               & Primärschlüssel, fortlaufende Nummer    \\
      \commandbox|ip|               & Die ersten drei Blöcke der IPv4-Adresse \\
      \commandbox|device|           & Desktop, mobile, tablet oder bot        \\
      \commandbox|brand_family|     & Z.B. Apple oder Samsung                 \\
      \commandbox|brand_model|      & Z.B. GALAXY S5                          \\
      \commandbox|browser_family|   & Z.B. Firefox, Chrome, Safari            \\
      \commandbox|browser_version|  & Die Versionsnummer des Browsers         \\
      \commandbox|platform_family|  & Z.B. Windows, Mac, Linux, iOS, Android  \\
      \commandbox|platform_version| & Die Versionsnummer der Plattform        \\
      \commandbox|user_id|          & Benutzer, der sich das Foto angesehen hat \\
      \commandbox|photo_id|         & Angesehenes Foto                        \\
      \commandbox|created_at|       & Zeitpunkt, als das Foto angesehen wurde \\
      \commandbox|updated_at|       & I.d.R. = \commandbox|created_at|        \\
    \end{tabular}
  \end{table}
\end{frame}

\begin{frame}[fragile]{Personalisierte Werbung}
  \begin{center}
    \Large
    \faCamera{} Beispiel auf \href{https://instahub.org}{InstaHub}
  \end{center}
\end{frame}

\begin{frame}{Auftrag}{Personalisierte Werbung}
  \textbf{Skript}
  \begin{itemize}
    \item \aufgabebox{Aufgabe 2.30}
    \item \challengebox{Aufgaben 2.31, 2.32}
  \end{itemize}
\end{frame}
